\section{Fehlerbetrachtung}

% \subsection{Temperaturabhängigkeit}

\subsection{Magnetfeldabhängigkeit}

Der Umrechnungsfaktor von Spulenspannung zu Magnetfeldstärke ($1.927\ \text{V} \rightarrow 5.8\ \text{mT}$) unterliegt Fertigungstoleranzen der Helmholtz-Spulen. Eine Abweichung verschiebt die Periodizität $\Delta B$ direkt proportional und führt zu einem systematischen Fehler in der Berechnung von $A_{\text{eff}}$ sowie folgend für $\lambda_L$ und $\lambda_J$.
An den Kontakten der Zuleitungen im Kryostaten entstehen durch Temperaturgradienten thermoelektrische Spannungen aufgrund des Seebeck-Effekts. Diese verschieben die gemessenen Spannungen geringfügig, was sich in der Auswertung durch den Plateau-Offset im Bereich von ca. $-0.09\ \text{mV}$ äußert. Da dieses Offset jedoch bei der Differenzbildung der kritischen Ströme rechnerisch eliminiert wurde, ist der Einfluss auf $I_0$ minimiert.
Der Vergleich der berechneten London'schen Eindringtiefe zu der aus \cite{kittel} zeigt eine hohe Abweichung. Unter Berücksichtigung der Probengröße, eventuellen Temperaturschwankungen und Oberflächeneffekten kann die Abweichung durchaus physikalisch sinnvoll sein und der Wert kann als passend im Rahmen der Messgenauigkeit angenommen werden.
Zusammenfassend zeigt die Übereinstimmung der Messdaten mit dem theoretischen Fraunhofer-Modell (relativer Fehler der Hauptparameter $I_0$ und $\Delta B$ von unter $6\ \%$), dass der Josephson-Kontakt im untersuchten Feldbereich makroskopisch korrekt arbeitet und die fundamentale Flussquantisierung im Supraleiter erfolgreich demonstriert.